% !PRINTED 408
\weberpage{408}
%<<CONT
und da $\sigma$ reell ist:

\[
\begin{aligned}
\Delta>0. &\qquad \sigma = \frac{1}{\sqrt{\Delta}}\overset{k}{\sum}(\Delta,k)\overset{n}{\sum}\frac{1}{n}\cos\frac{2nk\pi}{\Delta},\\
\Delta<0. &\qquad \sigma = \frac{1}{\sqrt{-\Delta}}\overset{k}{\sum}(\Delta,k)\overset{n}{\sum}\frac{1}{n}\sin\frac{mk\pi}{-\Delta}. % [sic] index \textquotedbl{}m\textquotedbl{} as printed, sum is taken over n
\end{aligned}
\]

Die beiden unendlichen Reihen nach $n$ haben aber bestimmte
Werte, die \gesperrt{Abel} in der Abhandlung über die Binomialreihe
aus der Potenzwickelung von $\log(1-z)$ abgeleitet hat, und die sich % [sic: Potenzwickelung as printed]
auch aus der Theorie der \gesperrt{Fourier}schen Reihen ergeben. Wird
dann $k$ positiv und zwischen $0$ und $\pm\Delta$ genommen, so ergibt sich:

\[
\begin{aligned}
\sum\frac{1}{n}\cos\frac{2nk\pi}{\Delta} &= -\log\left(2\sin\frac{k\pi}{\Delta}\right),\\
\sum\frac{1}{n}\sin\frac{2nk\pi}{-\Delta} &= \frac{\pi}{2}\left(1 - \frac{2k}{-\Delta}\right).
\end{aligned}
\]

Also ergibt sich aus (8) [mit Rücksicht auf $\Sigma(\Delta,k) = 0$]:

\weq{10}{\begin{aligned}
h\log\varepsilon &= -\overset{k}{\Sigma}(\Delta,k)\log\sin\frac{k\pi}{\Delta}, & \Delta&>0,\\
h &= \frac{1}{\Delta}\overset{k}{\Sigma}(\Delta,k)\,k, & \Delta&<0.
\end{aligned}}

Für den Fall der negativen Diskriminante kann man den
Ausdruck für $h$ so umgestalten, daß er die Form einer ganzen
Zahl annimmt.

Ist zunächst $\Delta$ ungerade ($\equiv 1 \bmod 4$), so kann man die
Zahlenreihe $k$ so zerlegen:

\[
\begin{aligned}
v, &\quad -\Delta - v,\\
2v, &\quad -\Delta - 2v,
\end{aligned}
\qquad 0 < v < -\frac{\Delta}{2},
\]

und erhält also, da hier $(\Delta,-v) = -(\Delta,v)$ ist (§ 85), zwei
Formen des Ausdrucks $h$:

\[
\begin{aligned}
h &= \frac{2}{\Delta}\overset{v}{\Sigma}(\Delta,v)\,v + \overset{v}{\Sigma}(\Delta,v),\\
h &= \frac{4}{\Delta}\Sigma(\Delta,2v)\,v + \overset{v}{\Sigma}(\Delta,v),
\end{aligned}
\]

woraus, wenn man die erste mit 2, die zweite mit $(\Delta,2)$ multipliziert
und subtrahiert:

\weq{11}{[2 - (\Delta,2)]\,h = \overset{v}{\Sigma}(\Delta,v).}

Wenn also $(\Delta,2) = -1$ ist, so muß die rechte Seite
durch 3 teilbar sein.
